\documentclass[11pt]{article}
\usepackage[T1]{fontenc}
\usepackage[utf8]{inputenc}
\usepackage{amsmath,amssymb,amsthm}
\usepackage{booktabs}
\usepackage{algorithm}
\usepackage{algpseudocode}
\usepackage{enumitem}
\usepackage{mathtools}
\usepackage[margin=1in]{geometry}
\usepackage{hyperref}

\renewcommand{\thesection}{S\arabic{section}}
\renewcommand{\thetable}{S\arabic{table}}
\renewcommand{\thealgorithm}{S\arabic{algorithm}}

\title{Supplementary Material\\
\large An Adaptive Large Neighborhood Search with Lagrangian Certification
for the Selective Routing Problem with Synchronization}
\author{}
\date{}

\begin{document}
\maketitle

\noindent
This document collects material referenced from the article. The main article
states the central model, validity proposition, and solver control flow; this
supplement provides the full arc-flow formulation, the synchronisation-aware
insertion and feasibility algorithms, the expanded Lagrangian derivation, and
the numerical safeguarding details needed for complete reproducibility.

\paragraph{Guide to this supplement.} Section~S1 gives the full SRPS-1 model.
Section~S2 collects Algorithms~S1--S3 (insertion, feasibility, and ejection
chains). Section~S3 gives the Lagrangian derivation, exact dynamic-programming
subproblem, and floating-point analysis; Section~S4 reports extended
computational evidence; and Section~S5 records the reproducibility and
certificate materials. These details extend, rather than replace, the main
article's definitions, validity argument, and headline evidence.

\section{Full arc-flow formulation}

We adopt the compact arc-flow SRPS-1 structure of Riera-Ledesma and Salazar-Gonz\'alez \cite{RieraLedesma2021SRPS}. Let $V_k=J_k\cup\{0,n+1\}$ be processor $k$'s node set (start depot $0$, end depot $n+1$). For each $k\in K$, let $z_{i\ell}^k\in\{0,1\}$ indicate whether processor $k$ traverses arc $(i,\ell)$ and let $x_{j,k}\in\{0,1\}$ indicate whether processor $k$ serves job $j$. Let $y_j\in\{0,1\}$ indicate whether job $j$ is selected. For timing, let $s_{j,k}$ be the service-start time of processor $k$ at job $j$ (for $j\in J_k$).

The objective is
\[
\max \sum_{j\in J} b_j y_j.
\]

Route structure is enforced by standard flow constraints:
\[
\sum_{\ell\in V_k\setminus\{0\}} z_{0\ell}^k = 1,\qquad
\sum_{i\in V_k\setminus\{n+1\}} z_{i,n+1}^k = 1,
\]
\[
\sum_{i\in V_k\setminus\{j\}} z_{ij}^k = x_{j,k},\qquad
\sum_{\ell\in V_k\setminus\{j\}} z_{j\ell}^k = x_{j,k}
\quad \forall j\in J_k,
\]
together with elementarity (no repeated job service on a route).

Joint-service coupling is
\[
y_j \le x_{j,k} \qquad \forall j\in J,\; k\in K_j,
\]
so a profit can be collected only if every required processor serves the job.

Timing and synchronisation are enforced by arc-time propagation and equal-start constraints. For each used arc $(i,\ell)$ on processor $k$,
\[
s_{\ell,k} \ge s_{i,k} + t_{i\ell} - M\bigl(1-z_{i\ell}^k\bigr),
\]
and for each selected job $j$ and each required pair $k,k'\in K_j$,
\[
s_{j,k}=s_{j,k'}.
\]
The global horizon is imposed by
\[
s_{n+1,k}-s_{0,k} \le L \qquad \forall k\in K.
\]

This explicit model resolves the synchronisation semantics: selected multi-processor jobs require both route inclusion by all processors in $K_j$ and simultaneous service start across those processors. Appendix A of the article keeps the bound-validity derivation and the exact form of the relaxed Lagrangian subproblems.

\section{Algorithmic details}

\subsection{Synchronisation-aware insertion}

The repair operators of the ALNS section of the article differ only in the order in which they offer candidate jobs to the insertion procedure below, which is where the synchronisation constraint is enforced.

\begin{algorithm}[t]
\caption{Synchronisation-aware insertion of a job (the repair primitive called within \textsc{ALNSPhase})}
\label{alg:insert}
\begin{algorithmic}[1]
\Require job $j$, current routes $\{P_k\}$, selected set $S$, instance \textit{inst}, per-processor candidate cap $c$
\Ensure the $m$ cheapest feasible cross-route placements of $j$, ranked by makespan increase ($m=1$ by default), or \textsc{Infeasible} if none exists
\State $R \gets K_j$ \Comment{processors that must jointly serve $j$}
\If{$R=\emptyset$} \Return \textsc{Infeasible} \EndIf
\State $e_0 \gets \textsc{ScheduleEnd}(\{P_k\}, S, \textit{inst})$ \Comment{makespan before inserting $j$}
\For{each processor $k \in R$}
    \State $C_k \gets [\,]$
    \For{each insertion position $p$ in route $P_k$}
        \State $\delta \gets t_{P_k[p-1],\,j} + t_{j,\,P_k[p]} - t_{P_k[p-1],\,P_k[p]}$ \Comment{local detour cost}
        \State append $(\delta,p)$ to $C_k$
    \EndFor
    \State $C_k \gets$ the $c$ entries of $C_k$ with smallest $\delta$
\EndFor
\State $\mathcal{F} \gets [\,]$ \Comment{feasible placements with their makespan increase}
\For{each combination $(p_k)_{k\in R} \in \prod_{k\in R} C_k$}
    \State $\{P'_k\} \gets \{P_k\}$ with $j$ inserted at position $p_k$ on every $k\in R$
    \State $(\textit{feasible},\, e) \gets \textsc{LongestPathScheduleCheck}(\{P'_k\}, S\cup\{j\}, \textit{inst})$
    \If{\textit{feasible}}
        \State append $\bigl(e-e_0,\ (p_k)_{k\in R}\bigr)$ to $\mathcal{F}$
    \EndIf
\EndFor
\State sort $\mathcal{F}$ by ascending $e-e_0$
\State \Return the first $m$ entries of $\mathcal{F}$ (\textsc{Infeasible} if $\mathcal{F}=[\,]$)
\end{algorithmic}
\end{algorithm}

\subsection{Synchronisation feasibility oracle}

Algorithm~\ref{alg:insert} calls \textsc{LongestPathScheduleCheck} as its feasibility primitive; this section gives that primitive in full, since it is where SRPS's synchronisation semantics are actually enforced and it is called on every candidate placement in Algorithm~\ref{alg:insert} and on every full-solution feasibility query elsewhere in the search.

The construction uses \emph{one node per job}, not one node per processor visit: if job $j$ is served jointly by processors $k$ and $k'\in K_j$, both routes contribute an incoming edge to the \emph{same} node $j$, and the longest-path relaxation below takes the maximum over all of a node's incoming edges. This is what enforces the common start time --- job $j$'s earliest start is the latest arrival among all processors required to serve it, with no separate cross-processor arc construction needed beyond the shared node.

\begin{algorithm}[t]
\caption{Synchronisation feasibility oracle (\textsc{LongestPathScheduleCheck})}
\label{alg:oracle}
\begin{algorithmic}[1]
\Require routes $\{P_k\}_{k\in K}$, selected set $S$, instance \textit{inst} (transition times $T$, horizon $L$, depot nodes $\textit{start}, \textit{end}$)
\Ensure $(\textit{dist}, \textit{feasible})$ --- earliest start times and a feasibility verdict
\State $V \gets \{\textit{start}, \textit{end}\} \cup S$ \Comment{one node per job; no per-processor duplication}
\State $E \gets \emptyset$
\For{each route $P_k$}
    \For{each consecutive pair $(i,\,i')$ in $P_k$}
        \If{$(i,i') \notin E$} \Comment{dedupe: identical arc may recur across processors}
            \State add arc $i \to i'$ with weight $T[i][i']$ to $E$
        \EndIf
    \EndFor
\EndFor
\State $(\textit{topo}, \textit{acyclic}) \gets \textsc{KahnTopologicalSort}(V, E)$ \Comment{cycle detection}
\If{$\lnot \textit{acyclic}$}
    \State \Return $(\textit{None}, \textsc{Infeasible})$ \Comment{contradictory synchronisation ordering}
\EndIf
\State $\textit{dist}[\textit{start}] \gets 0$; \ $\textit{dist}[v] \gets -\infty$ for all other $v \in V$
\For{each $u \in \textit{topo}$ in topological order} \Comment{single relaxation pass}
    \For{each arc $(u, v, w) \in E$}
        \State $\textit{dist}[v] \gets \max\bigl(\textit{dist}[v],\ \textit{dist}[u] + w\bigr)$ \Comment{max over ALL incoming arcs to $v$ --- this is the synchronisation}
    \EndFor
\EndFor
\If{$\textit{dist}[\textit{end}] > L$}
    \State \Return $(\textit{dist}, \textsc{Infeasible})$ \Comment{horizon violated}
\Else
    \State \Return $(\textit{dist}, \textsc{Feasible})$
\EndIf
\end{algorithmic}
\end{algorithm}

Correctness rests on two conditions, checked in this order: (i) $E$ induces no directed cycle on $V$ --- a cycle would require some processor to serve a job before its own predecessor, which Kahn's algorithm detects as a subset of $V$ left with nonzero in-degree after the queue empties; (ii) the longest path from $\textit{start}$ to $\textit{end}$, which is exactly $\textit{dist}[\textit{end}]$ after the relaxation pass, does not exceed $L$. Because $E$ is acyclic whenever the check reaches the relaxation pass, a single topological-order sweep computes every $\textit{dist}[v]$ exactly (no re-relaxation is needed, unlike Bellman--Ford on a general graph). With $|V| = O(|S|)$ and $|E| = O\bigl(\sum_{j\in S}|K_j|\bigr) = O(|S|)$ since $|K_j|\in\{2,3,4\}$ is bounded, both the topological sort and the relaxation pass run in $O(|S|)$ time; the oracle is called once per candidate placement in Algorithm~\ref{alg:insert} and once per full-solution feasibility query.

\emph{Worked example.} Job $j$ is served by processors $k$ and $k'$. On $k$'s route, $j$'s predecessor is $a$, giving arc $(a, j)$ with weight $T_{a,j}$; on $k'$'s route, $j$'s predecessor is $b$, giving arc $(b, j)$ with weight $T_{b,j}$. Both arcs point into the same node $j$. If $\textit{dist}[a] + T_{a,j} = 40$ and $\textit{dist}[b] + T_{b,j} = 55$, the relaxation sets $\textit{dist}[j] = 55$: job $j$ cannot start until the later-arriving processor $k'$ is ready, and the $15$ units of idle time this imposes on $k$'s route are absorbed automatically by $k$'s own subsequent arc weights, without any explicit ``waiting'' variable or synchronisation arc distinct from the two ordinary route arcs already present.

\subsection{Ejection-chain post-processing}

Algorithm 1 of the article calls \textsc{EjectionChain} once after the destroy--repair loop terminates. It targets a specific failure mode: a job $j$ that the insertion procedure of Algorithm~\ref{alg:insert} could not place because every candidate position it tried was blocked by an already-served job, even though evicting that job and re-inserting it elsewhere would free the position at no net cost.

\begin{algorithm}[t]
\caption{Ejection-chain post-processing}
\label{alg:ejection}
\begin{algorithmic}[1]
\Require solution $(\textit{selected}, \{P_k\})$, max depth $d\in\{1,2\}$, max rounds $R$
\Ensure improved $(\textit{selected}, \{P_k\})$ in place; total profit gained
\State $\textit{gained} \gets 0$
\For{$\textit{round} = 1$ \textbf{to} $R$}
    \State $U \gets$ unserved jobs sorted by profit, descending
    \State $\textit{improved} \gets \textsc{False}$
    \For{each $j \in U$} \Comment{(a) direct insertion sweep}
        \If{$j$ has a feasible placement (Algorithm~\ref{alg:insert})}
            \State insert $j$; $\textit{gained} \gets \textit{gained} + b_j$; $\textit{improved} \gets \textsc{True}$; \textbf{break} to next round
        \EndIf
    \EndFor
    \If{$\textit{improved}$} \textbf{continue} \EndIf
    \For{each $j \in U$} \Comment{(b) eject 1 (or, if $d=2$, up to 2) job(s), then insert $j$}
        \State $C \gets$ jobs occupying $j$'s required-processor routes, sorted by profit ascending \Comment{depth 2: cheapest 8 only}
        \For{each candidate eviction set $Q\subseteq C$, $|Q|\le d$}
            \State $\textit{trial} \gets$ copy of the solution with $Q$ removed
            \If{$j$ has a feasible placement in $\textit{trial}$ (Algorithm~\ref{alg:insert})}
                \State insert $j$ into $\textit{trial}$
                \State try to re-insert every $q\in Q$ elsewhere in $\textit{trial}$ (Algorithm~\ref{alg:insert})
                \State $\Delta \gets b_j - \sum_{q\in Q \text{ not re-inserted}} b_q$
                \If{$\Delta > 0$}
                    \State commit $\textit{trial}$; $\textit{gained} \gets \textit{gained}+\Delta$; $\textit{improved}\gets\textsc{True}$; \textbf{break} to next round
                \EndIf
            \EndIf
        \EndFor
    \EndFor
    \If{$\lnot\textit{improved}$} \textbf{break} \EndIf
\EndFor
\State \Return $\textit{gained}$
\end{algorithmic}
\end{algorithm}

Every successful move restarts the sweep over unserved jobs from the highest-profit end, since a single eviction can free positions for jobs other than the one it was performed for. At depth~2, the eviction-candidate set $C$ is capped at its eight cheapest members before forming pairs, keeping $\binom{|C|}{2}$ tractable; the article's Algorithm 1 uses depth~3 in production, which follows the identical pattern with $|Q|\le 3$ and a cap of five candidates. The pass terminates when a full round finds no improving move, or after $R=30$ rounds.

\section{Lagrangian-bound derivation, exact subproblem solution, and numerical safeguarding}
\label{sec:supp-lagrangian}

This section expands the Lagrangian-bound material referenced from the article's
Section~3: the full algebraic derivation of the per-processor decomposition, the
bitmask dynamic-programming recurrence, and the floating-point safeguarding
analysis behind the certified gap. The article states the model, the relaxed
bound, and the central validity proposition (with proof) in full; nothing here
is required to assess correctness or novelty, and this section is provided so
the numerical treatment can be reconstructed and checked independently.

\subsection{Expanded derivation of the per-processor decomposition}

We dualise the joint-service coupling $y_j\le x_{j,k}$ ($j\in J,\ k\in K_j$) with
multipliers $\mu_{j,k}\ge 0$. For $\mu\ge 0$,
\[
L(\mu) = \max_{(y,\{P_k\})}\ \sum_{j\in J} b_j y_j + \sum_{j\in J}\sum_{k\in K_j}\mu_{j,k}\,(x_{j,k}-y_j),
\]
the maximisation being over $y\in\{0,1\}^n$ and independent per-processor
elementary routes within horizon $L$, with the inter-processor synchronisation
requirement relaxed (each processor's route is otherwise still constrained to be
elementary and within $L$). Regrouping the objective by the variables it
multiplies,
\[
\sum_{j\in J} b_j y_j + \sum_{j,k}\mu_{j,k}(x_{j,k}-y_j)
= \sum_{j\in J}\Bigl(b_j-\!\sum_{k\in K_j}\mu_{j,k}\Bigr)y_j \;+\; \sum_{k\in K}\sum_{j\in J_k}\mu_{j,k}x_{j,k},
\]
and because $y_j$ now appears only in the first sum and is otherwise
unconstrained, the maximising choice is $y_j=1$ exactly when its coefficient is
positive, giving $\max(0,\ b_j-\sum_{k\in K_j}\mu_{j,k})$ for that term. The
second sum decomposes over processors because each $x_{j,k}$ is constrained only
by processor $k$'s own routing feasibility, giving
\[
L(\mu) = \sum_{j\in J}\max\!\Bigl(0,\; b_j-\!\sum_{k\in K_j}\mu_{j,k}\Bigr) \;+\; \sum_{k\in K} O_k(\mu),
\qquad
O_k(\mu) = \max\Bigl\{ \textstyle\sum_{j\in S}\mu_{j,k} : S\subseteq J_k,\ S\text{ admits an elementary route within }L \Bigr\}.
\]
$O_k(\mu)$ is a single-processor orienteering problem with item profits
$\mu_{j,k}$: select a subset of jobs and an elementary route serving them within
the horizon, maximising the collected multiplier value. This is the
decomposition the article's Proposition~1 certifies as valid for every
$\mu\ge 0$; both operations used --- dualising the coupling and relaxing the
synchronisation requirement inside the subproblem --- only enlarge the feasible
region or add a non-negative penalty term, which is exactly what the proof in
the main text uses.

\subsection{Exact solution of $O_k(\mu)$: bitmask recurrence and cost}

Jobs with $\mu_{j,k}\le 0$ never improve an orienteering objective and are
removed a priori. Writing $\mathrm{dp}[\textit{mask}][i]$ for the minimum time to
serve exactly the job set $\textit{mask}\subseteq J_k$ and end at job $i\in
\textit{mask}$, the recurrence is
\[
\mathrm{dp}[\{i\}][i] = t_{0,i}, \qquad
\mathrm{dp}[\textit{mask}][i] = \min_{i'\in \textit{mask}\setminus\{i\}} \mathrm{dp}[\textit{mask}\setminus\{i\}][i'] + t_{i',i},
\]
and $O_k(\mu)$ is the largest profit $\sum_{j\in\textit{mask}}\mu_{j,k}$ over
masks whose best completion back to the depot, $\min_i \mathrm{dp}[\textit{mask}][i] + t_{i,0}$,
is within $L$. This enumerates the optimal route over all $2^{|J_k|}$ subsets and
$|J_k|$ endpoints and is therefore \emph{exact}: no pruning, sampling, or
candidate-set restriction is used. Its cost is $O(2^{|J_k|}|J_k|^2)$.

This is tractable at every benchmark scale because $|J_k|$ is small and does not
grow with $n$: each job requires only $|K_j|\in\{2,3,4\}$ of the $|K|=55$
processors, so the incidences spread thinly. The per-instance mean $|J_k|$ is
$2.8$--$6.0$ and the maximum over all $660$ study instances is $13$, attained
already at $n=80$ and unchanged at $n=150$; every subproblem therefore costs at
most $2^{13}\cdot 13^2\approx 1.4\times10^6$ operations.

\subsection{Floating-point safeguarding of the reported certificate}
\label{sec:supp-numerical}

Every iterate $L(\mu_t)$ is a valid upper bound by the article's Proposition~1,
so $\min_t L(\mu_t)$ is valid; since all profits $b_j$ are integral, the article
reports $UB_{\text{lag}}=\lfloor\min_t L(\mu_t)\rfloor$, which is also valid
because no feasible objective lies strictly between the two. One numerical point
deserves care, because it bears directly on validity rather than tightness. The
argument above is exact in real arithmetic, but $L(\mu)$ is evaluated in
floating point as a sum of many terms, and its computed value can fall
marginally below the exact one. Flooring then removes an entire profit unit and
the resulting quantity is no longer a valid upper bound. The situation is not
pathological: as the subgradient converges, $L(\mu)$ approaches the integral
optimum \emph{from above}, so it sits adjacent to an integer precisely when
convergence is attained --- exactly on the instances a certificate is most
likely to close. Obtaining bounds that remain valid under floating-point
evaluation requires an explicit numerical safeguard, and we adopt the remedy: the floor is
taken with a tolerance,
\[
UB_{\text{lag}} = \bigl\lfloor \min_t L(\mu_t) + \varepsilon \bigr\rfloor,
\qquad \varepsilon = 10^{-9},
\]
far below one profit unit and therefore unable to weaken the bound, while
removing the failure mode entirely. A reported bound is additionally never
allowed below the incumbent objective, since a feasible solution attains that
value; with the tolerance applied this guard is never triggered.

This safeguard is not hypothetical. Computed without the tolerance, the defect
materialises on four of the $660$ study instances
(\texttt{C\_n080\_034\_a75\_102}, \texttt{D\_n065\_028\_a75\_084},
\texttt{D\_n070\_035\_a75\_105} and \texttt{ED\_n045\_008\_a75\_024}): on each,
the computed bound falls exactly one profit unit below its exact value and
thereby coincides with the incumbent, certifying a false optimality. All four
lie at the loosest budget level ($\alpha=0.75$), which is what the mechanism
predicts --- subgradient convergence drives $L(\mu)$ onto an integer \emph{from
above}, so the flooring error is most exposed precisely where the dual has
converged best. With the tolerance in place the four bounds are each corrected
upward by one unit and the instances are reported with certified gaps of
$0.029\%$, $0.047\%$, $0.040\%$ and $0.068\%$ rather than as proven optimal.
Every certificate reported in the article is produced with the tolerance
applied.

\subsection{Independent re-derivation (Path~B) and its result}

A rigorously safe bound calls for directed rounding or interval arithmetic to
bound total floating-point error
across all arithmetic operations, independently of the search that produced the
certificate. Algorithm~\ref{alg:pathb} re-derives $L(\mu)$ for every one of the
$660$ certificates from two inputs only: the raw benchmark file and the
certificate's stored multipliers $\mu$, with no dependence on any cached or
in-process production object. Every float64 input is converted to its exact
decimal value ($\texttt{Decimal(x)}$, which reproduces the input's binary value
bit-for-bit, not the lossy string round-trip $\texttt{Decimal(str(x))}$), and
every summation is carried out in $80$-digit decimal arithmetic under
round-toward-$+\infty$ (\texttt{ROUND\_CEILING}), giving a directed-rounding
upper bound on $L(\mu)$ rather than a repetition of the float64 computation. The
per-processor orienteering subproblem is re-solved to obtain its selected job
subset, and that subset's profit total is independently re-summed in exact
arithmetic and checked against the value the bitmask DP reported.

\begin{algorithm}[t]
\caption{Independent certificate re-derivation (Path~B)}
\label{alg:pathb}
\begin{algorithmic}[1]
\Require certificate $c=(\textit{name},\mu,UB_c,\varepsilon)$, reloaded independently of the search
\Ensure verdict $\in\{\textsc{Pass},\textsc{PassEps},\textsc{Fail}\}$
\State $\textit{inst}\gets\textsc{LoadRawInstance}(c.\textit{name})$ \Comment{from the benchmark file, not a cached object}
\State $y_{\text{contrib}}\gets 0$ \Comment{exact decimal arithmetic, round-toward-$+\infty$}
\For{$j\in J$}
    \State $r_j\gets b_j-\sum_{k\in K_j}\mu_{j,k}$
    \If{$r_j>0$}\ $y_{\text{contrib}}\gets y_{\text{contrib}}+r_j$\ \EndIf
\EndFor
\State $\textit{total\_ok}\gets 0$
\For{$k\in K$}
    \State $S_k\gets\{j\in J_k:\mu_{j,k}>0\}$
    \State $(v_k,\textit{selected}_k)\gets\textsc{OrienteeringDP}(S_k,\mu_{\cdot,k},L)$ \Comment{exact bitmask DP, Section~\ref{sec:supp-lagrangian}}
    \State $\hat v_k\gets\sum_{j\in\textit{selected}_k}\mu_{j,k}$ \Comment{exact re-sum of the DP's own selected subset}
    \State \textbf{flag} if $|\hat v_k-v_k|>10^{-9}$ \Comment{DP-vs-exact discrepancy; none observed on any of $660\times 55$}
    \State $\textit{total\_ok}\gets\textit{total\_ok}+\hat v_k$
\EndFor
\State $L(\mu)\gets y_{\text{contrib}}+\textit{total\_ok}$
\If{$\lfloor L(\mu)\rfloor\ge UB_c$}\ \Return \textsc{Pass}
\ElsIf{$\lfloor L(\mu)+\varepsilon\rfloor\ge UB_c$}\ \Return \textsc{PassEps}
\Else\ \Return \textsc{Fail}
\EndIf
\end{algorithmic}
\end{algorithm}

The result: all $660$ certificates are valid at \emph{zero} tolerance under this
independent re-derivation --- none require $\varepsilon$ to floor to the correct
integer, and the per-processor exact re-sum matches the DP's own reported total
on every one of the $660\times 55$ subproblems checked (maximum observed
discrepancy $0.0$). This isolates the source of the four-instance failure mode
above precisely to float64's summation rounding, accumulated across the
$\sum_k|J_k|$ and $|J|$ terms of $y_{\text{contrib}}$ and $\textit{total\_ok}$,
and nothing else --- the DP's combinatorial subset selection is exact, and the
underlying mathematical value of $L(\mu)$, computed from the same stored
multipliers, was already a valid upper bound before any tolerance was applied.
The $\varepsilon=10^{-9}$ safeguard is therefore confirmed to compensate for a
real but purely arithmetic artefact, not to mask a genuine risk of invalidity,
on every instance examined. The verification script and its full per-instance
output for all $660$ certificates are included in the code and results
accompanying the article, and can be reproduced with a single command
(\texttt{reproduce.py --stage pathb}). The claim remains scoped to the studied
instances and this implementation: it is a direct, complete check of the
reported results, not a universal guarantee for arbitrary floating-point
environments or not-yet-examined instance classes.

\section{Extended computational results}
\label{sec:supp-extended}

This section gives the material referenced from the article's ablation,
sensitivity, and operating-point discussion in full: the article reports a
single pooled ablation table and headline sensitivity conclusions; the
per-family and per-size breakdowns, the full sensitivity sweep, the refinement
iteration-cap sweep, and the replayed early-stopping comparison are given here.

\subsection{Ablation: beneath the pooled mean}

The pooled means the article reports can mask cancellation: a component that
helps on some instances and hurts on others can pool to near zero without being
inert anywhere. Counting instances individually against the measured noise
floor (mean $|\Delta|=0.036$\,pp, Tier~2, 70 hard-tail instances) rather than
averaging them: B4 exceeds it on $31/40$ instances and B2 on $7/40$, consistent
with their pooled means. B1 exceeds it on $11/40$ instances (reaching
$+0.47$\,pp on one), B3 on $8/40$ (reaching $+0.32$\,pp on one, $-0.16$\,pp on
another), and B6 on $5/40$ --- all three pool to near-zero only because their
per-instance effects run in both directions. B5, B7 and B8 remain null by this
finer count too ($2/40$, $1/40$ and $0/40$ respectively), and B7 additionally
changes the search's stop reason on \emph{zero} of the $40$ instances.

Per-family cells in this sample range from $1$ to $8$ instances, too few to
test formally; we report the following pattern as exploratory, not
confirmatory. B1's and, more markedly, B2's effect concentrates on families D
and ED (B2: $+0.22$\,pp on hard-tail D, $n=6$; $+0.12$\,pp on hard-tail ED,
$n=5$) and is indistinguishable from zero on families B, EB and EC in the same
stratum. Both arms' effects also scale with instance size, roughly doubling on
the larger half of the sample for every component measured. Neither pattern is
claimed beyond a hypothesis worth testing at larger $n$.

\subsection{Ablation: B5/B7 cross-check against the census and identity}

B5 (no ejection chain) and B7 (no in-loop re-bound) each have a stronger,
larger-population measurement already in hand (Section~6 of the article and the
ejection-chain identity), and are run as sampled arms here purely as a
cross-check, not as a replacement for those measurements. B5's sampled mean is
$+0.00$\,pp hard tail / $+0.01$\,pp closing, against the identity's $0.000$\,pp
/ $0.0011$\,pp over the full $660$ --- the same order of magnitude, and both
effectively zero on the hard tail, where the mechanism cannot fire before the
search has closed to within a few tenths of a percent. B7 is a cleaner check:
the returned objective is identical to A0 on $20/20$ sampled hard-tail instances
and $19/20$ closing instances, matching the census's $70/70$ exact tie on the
hard tail (Table~6 of the article). The re-bound's independently measured
$0.264$\,pp mean effect is a certificate effect invisible to the fixed-common-bound
ablation by construction --- the sampled arm did not contradict that argument,
it reproduced it.

\subsection{Full one-at-a-time sensitivity sweep}

Table~\ref{tab:sensitivity} reports the full sweep on the 30-instance stratified
subset; the baseline mean certified gap is $0.1525\%$, measured before
refinement for the same reason as the hardness analysis. A duplicate trial of
the tightest gap-threshold setting ($0.1\%$) measures the run-to-run noise floor
at $0.005$\,pp on this subset; nine of the eleven variants fall within that
resolution.

These gap figures are also on the pre-monotonicity-guard dual-bound codepath:
both the refinement stage and the \texttt{min(ub, new\_ub)} in-loop guard
(Section~\ref{sec:experiments}) were added to the production driver after
these trials ran, and \texttt{run\_sensitivity.py} invokes neither. Neither
change can affect this table's conclusions through the incumbent objective,
since refinement re-solves the dual against a fixed incumbent and never
alters it, and the primal search code was otherwise unchanged over this
period. As a refinement-independent check, the last column reports each
variant's mean objective deviation from the production reference incumbent,
computed from the same trial runs. Ten of the eleven variants reproduce the
reference within $0.01\%$, below a $0.0015$\,pp duplicate-run noise floor for
this metric; only the $450$\,s phase-budget variant shows a clear effect
($+0.011\%$), consistent with it being one of the two variants that exceed the
gap-based noise floor below. The two dual-only parameters, \texttt{lag\_max\_iter}
and \texttt{lag\_max\_time}, show exactly zero objective deviation from
baseline, as expected since they act only on the subgradient re-bound and
never touch the primal search.

\begin{table}[h]
\centering
\caption{One-at-a-time parameter sensitivity (mean certified gap, 30-instance subset). Obj.\ $\Delta$ is the refinement-independent mean objective deviation from the production reference incumbent, relative to the baseline trial.}
\label{tab:sensitivity}
\begin{tabular}{llccc}
\toprule
Parameter & Setting & Mean cert \% & $\Delta$ baseline (pp) & Obj.\ $\Delta$ vs ref.\ (pp) \\
\midrule
\textit{(baseline)} & 300\,s / 0.3\% / 25--33--40\% / 200 / 60\,s & 0.1525 & --- & 0.000 \\
phase\_rt        & 150\,s  & 0.1485 & $-0.004$ & $0.000$ \\
phase\_rt        & 450\,s  & 0.1390 & $-0.014$ & $+0.011$ \\
gap\_threshold   & 0.1\%   & 0.1261 & $-0.026$ & $+0.005$ \\
gap\_threshold   & 0.5\%   & 0.1543 & $+0.002$ & $-0.000$ \\
gap\_threshold   & 1.0\%   & 0.1505 & $-0.002$ & $-0.001$ \\
destroy\_fracs   & tight   & 0.1485 & $-0.004$ & $+0.005$ \\
destroy\_fracs   & wide    & 0.1487 & $-0.004$ & $+0.003$ \\
lag\_max\_iter   & 100     & 0.1570 & $+0.005$ & $0.000$ \\
lag\_max\_iter   & 400     & 0.1510 & $-0.002$ & $0.000$ \\
lag\_max\_time   & 30\,s   & 0.1566 & $+0.004$ & $0.000$ \\
lag\_max\_time   & 120\,s  & 0.1485 & $-0.004$ & $0.000$ \\
\bottomrule
\end{tabular}
\end{table}

Two variants exceed the noise floor: tightening the stopping threshold to
$0.1\%$ improves the mean by $0.026$\,pp, and extending the phase budget to
$450$\,s improves it by $0.014$\,pp. Both move in the expected direction and
neither was adopted in production, since both push compute toward the
one-hour per-instance cap by forcing early-closing instances to keep searching
and by extending each phase. The production defaults are a deliberate budget
choice, not an inadvertent local optimum.

\subsection{Refinement iteration-cap sensitivity}

Table~\ref{tab:refinesens} reports the refinement stage's subgradient
iteration-cap sweep, run on the $60$ loosest instances rather than the
stratified subset, since those are the only instances on which the cap can
bind; the mean gaps are correspondingly far above the $660$-instance headline
figures and are not comparable with them.

\begin{table}[h]
\centering
\caption{Refinement iteration-cap sensitivity, measured on the $60$ loosest
instances. Time is the median per-instance refinement cost.}
\label{tab:refinesens}
\begin{tabular}{lcccc}
\toprule
Iteration cap & Mean gap \% & Median gap \% & Closed & Median time \\
\midrule
$1000$            & 0.5061 & 0.3651 & 4/60  & 35\,s \\
$3000$ (used)     & 0.4666 & 0.3488 & 10/60 & 96\,s \\
$5000$            & 0.4643 & 0.3488 & 10/60 & 154\,s \\
\bottomrule
\end{tabular}
\end{table}

The subgradient satisfies its convergence test on $56$ of the $660$ instances;
on the rest it is still descending when the cap stops it, so every reported
bound is a valid but truncated dual iterate rather than a dual optimum.
Reducing the cap to $1000$ is measurably worse (six fewer instances closed, mean
up $0.040$\,pp); raising it to $5000$ buys $0.002$\,pp and no additional
closures for roughly $60\%$ more time. The chosen value of $3000$ sits at the
knee of the curve: any value from a few thousand upward gives the same
certificates.

\subsection{Replayed early-stopping operating points}

Table~\ref{tab:operating} answers what the search would give up if it stopped
earlier, by replaying two simple stall-based stopping rules against the
recorded phase trajectories rather than re-running the solver. Replaying is
exact on the primal side because the stall tolerance enters only the
termination test: a run stopped at the $\kappa$-th consecutive stall is a
strict prefix of a run stopped later. This was checked directly: on four
hard-tail instances the state read from the trajectory matched a separate
execution configured to stop at the first stall exactly (objectives identical,
runtimes within $0.6\%$), and on a 30-instance sample instrumented to log the
solver's own stall counter, the replayed stopping phase coincided with it on
all $30$ instances at each of $\kappa=1,2,3$, with certified gaps at each
stopping point reproducing the replayed values to within $0.005$\,pp on
average.

\begin{table}[h]
\centering
\caption{Operating points, over the $660$-instance study set. Each row is the
recorded state at the stated stopping condition, with refinement applied. The
first row is the configuration used throughout the article.}
\label{tab:operating}
\begin{tabular}{lccccccccc}
\toprule
& \multicolumn{5}{c}{Certified gap} & \multicolumn{3}{c}{Runtime (min)} & $z$ lost \\
\cmidrule(lr){2-6}\cmidrule(lr){7-9}
Stopping rule & mean \% & median \% & max \% & proven opt. & $\le 0.5\%$ & median & mean & max & (inst./units) \\
\midrule
Full search (used) & 0.0658 & 0.0000 & 1.6135 & 446/660 & 96.8\% & 6.2 & 8.6 & 46.9 & --- \\
Stop at 2nd stall  & 0.0663 & 0.0000 & 1.6135 & 446/660 & 96.8\% & 6.2 & 7.7 & 46.9 & 4 / 11 \\
Stop at 1st stall  & 0.0683 & 0.0000 & 1.6135 & 443/660 & 96.7\% & 5.4 & 6.0 & 30.7 & 12 / 43 \\
\bottomrule
\end{tabular}
\end{table}

The trade is mild. Stopping at the first stall costs $0.0026$\,pp of mean
certified gap and three of the $446$ optimality proofs, while reducing mean
runtime by $30\%$ and worst-case runtime by $35\%$; the median runtime barely
moves, since instances certified at construction finish in under a second
whatever the stopping rule and anchor the median. The saving is concentrated in
the tail, where it is worth having: the primal cost is $43$ profit units in
total across all $660$ instances (largest single loss $13$), and the reason it
is mild is refinement --- the certificate no longer depends on how long the
search ran, so truncating the search costs only what the primal loses, and the
primal has usually stopped improving anyway.

These are reported as operating points, not a recommendation, and neither is
adopted in production. Two limitations apply: the rules are replayed from a
single recorded run per instance, so the frontier describes the trade-off this
run realised rather than an expectation over runs; and the refined bounds were
computed against the final incumbent, so on instances where truncation lowers
$z$ the corresponding gaps are marginally optimistic.

\section{Reproducibility and certificates}

The accompanying repository contains the source code, analysis scripts,
versioned result CSVs, and the canonical 660-instance multiplier certificates
under \texttt{results/certificates/}. Benchmark inputs are obtained separately
from the public OPS suite cited in the article. The reproducibility map in
\texttt{docs/REPRODUCIBILITY.md} identifies the input artifacts and command for
every numerical manuscript result.

For a deterministic rebuild of current tables, figures, and certificate checks,
run \texttt{python reproduce.py --manuscript-results}. The independent
certificate check is \texttt{python verify\_interval\_arithmetic.py --mode exact
--csv}; it reloads each raw instance and stored multiplier vector, then
recomputes the reported Lagrangian value with directed rounding. The archived
primary run and post-search refinement are distinct versioned series, as the
latter changes certificates but never incumbents. Time-budgeted campaign reruns
may differ in wall-clock iteration counts and are documented separately in the
same map. Historical companion-workspace and licensed CPLEX drivers are not
required by the standalone manuscript-results command.

\bibliographystyle{plain}
\bibliography{refs}

\end{document}
